% Nouveau EGA V (cote 25): archive cote 25, site page 167 (PDF page 168).
% typescript. Draft compared: none.
% AI-transcribed by Claude Opus 5.5 (claude-opus-5-5, adaptive thinking, effort high), Message Batches API, batch ega25b (msgbatch_01A7n5YuXubUPNK1zsepi8RG), 2026-11-13; original AI transcription, subsequently checked against both supplied scan parts; see PAGES_CHECKED.tsv and CHECK_DECISIONS.tsv.
\scanpage{25/167}
\mspage{\hand{2}}

\[ v^1 : E\hand{\otimes}\underline{O}_{V^1} \hand{\longrightarrow} \hand{\underline{\Omega}}^1_{X/Y}\hand{\otimes}\underline{O}_{V^1} \]
\var{\xout{en termes duquel}}{}{typewriter strike-out}{ega25_25_167-v0001} (déduit  lui-même de l'homomorphisme de modules
localement libres sur $V^o$
\[ v : \underline{E}\hand{\otimes}\underline{O}_{V^o} \hand{\longrightarrow} \hand{\underline{\Omega}}^1_{X/Y}\hand{\otimes} \underline{O}_{V^o} \quad , \]
en termes duquel  $V^1$ a été défini comme \var{\inserted[typed]{le}}{le}{insertion}{ega25_25_167-v0002} sous-préschéma de $V^o$ sur lequel
le rang de v est strictement \hand{$\leqslant$} m-1 = rang $E\hand{\otimes}\underline{O}_{V^o}$ - 1). Pour la définition
du premier isomorphisme, cf page 5, pour celle du second, df compléments
diffétentiels (ci joints) à `` variation du rang d'un homomorphisme de
Modules loc libres''. On obtient donc une suite exacte:
\[ 0 \hand{\to} \underline{E}\hand{\otimes}\underline{O}_{V^1} \hand{\xrightarrow{\ \textcolor{black}{i}\ }} \underline{N}_{V^1/X} \hand{\xrightarrow{\ \ \textcolor{black}{j}\ \ }} \underline{\mathrm{Hom}}_{\underline{O}_{V^1}}(W,K) \hand{\longrightarrow} 0 \qquad . \]
Il faut alors vérifier \var{\deleted{\var{\xout{1\textsuperscript{o}) que}}{}{typewriter strike-out}{ega25_25_167-v0004}}}{}{deletion}{ega25_25_167-v0003}
1\textsuperscript{o}) que
\[ \hand{\mu}i : \underline{E}\hand{\otimes}\underline{O}_{V^1} \hand{\longrightarrow} \hand{\underline{\Omega}}^1_{X/Y}\hand{\otimes}\underline{O}_{V^1} \]
n'est autre que $v^1$ , ce qui implique alors , grâce à l'hypothèse que
$v^1$ est strictement de rang m-1, que \var{\xout{V(u)}}{}{typewriter strike-out}{ega25_25_167-v0005} $\hand{\mu}i$ définit par passage au
quotient
\[ \hand{/\mu'} : \quad \underline{\mathrm{Hom}}\,(W,K) \hand{\longrightarrow} W = \mathrm{Coker}\ v^1 \ , \]
et que V(u) est aussi le sous-préschéma de $V^1$ qui rend $\hand{\mu'}$ strictement
de rang \hand{\unsure{$<|\leqslant$}}rang W \var{\deleted{\unsure{$-$}1}}{}{deletion}{ega25_25_167-v0006} ; d'autre part, il' \var{\xout{peut aussi s'interpréter}}{}{typewriter strike-out}{ega25_25_167-v0007} faut véri-%
fier

2\textsuperscript{o}) que $\hand{\mu'}$ n'est autre que l'homomorphisme déduit de la forme
quadratique \var{\xout{sur}}{}{typewriter strike-out}{ega25_25_167-v0008} M sur W à valeurs dans K, du bas de page 15. Il en résulte
aussitôt que V($\hand{\mu'}$) = V(\unsure{D}) \var{\xout{x}}{}{typewriter strike-out}{ega25_25_167-v0009}, où D est le discriminant de M, ce \var{\xout{qui\unsure{t|l}e}}{}{typewriter strike-out}{ega25_25_167-v0010} qui
achève la démonstration.

La vérification de 1\textsuperscript{o}) et 2\textsuperscript{o}) (où peut-être il y a à introduire un
signe ?) ne peut être que essentiellement méchanique, et est laissée à
Dieudonné.

\mspage{\hand{167}}
